\documentclass[10pt,aspectratio=169,serif]{beamer}

\usepackage[fontset=none, scheme=plain]{ctex}
\IfFontExistsTF{Songti SC}{%
	\setCJKmainfont{Songti SC}[AutoFakeSlant=0.2]%
	\setCJKsansfont{Songti SC}[AutoFakeSlant=0.2]%
}{%
	\setCJKmainfont{FandolSong-Regular}[Extension=.otf,BoldFont=FandolSong-Bold,AutoFakeSlant=0.2]%
	\setCJKsansfont{FandolSong-Regular}[Extension=.otf,BoldFont=FandolSong-Bold]%
}
\usepackage{amsmath,amssymb,graphicx,tikz,booktabs,multicol,calligra}
\usepackage{pgfplots}
\pgfplotsset{compat=1.18}
\usepackage{tcolorbox}
\tcbuselibrary{skins}
\usepackage{XyygU}
\usepackage{caption}
\captionsetup[figure]{name=图,labelsep=period}

\usetikzlibrary{arrows.meta,positioning,calc}

\setlength{\abovedisplayskip}{0pt plus 2pt minus 2pt}
\setlength{\belowdisplayskip}{0pt plus 2pt minus 2pt}

\title[习题2.1]{习题2.1 等式性质与不等式性质}
\subtitle{复习巩固 \& 综合运用 \& 拓广探索}
\author{}
\institute{}
\date{2026-10-16}

\begin{document}

\begin{frame}
  \titlepage
\end{frame}

% ========== 第1题 ==========
\begin{frame}[t]{第1题：生活中的不等关系}
\textbf{1.} 举出几个现实生活中与不等式有关的例子。
\end{frame}

\begin{frame}[t]{第1题 解答}
\begin{tcolorbox}[colback=orange!10,colframe=orange,title=\textbf{参考答案}]
\textbf{举例：}
\begin{itemize}
  \item 限速 $v \le 120\ \mathrm{km/h}$；
  \item 购票身高 $h \ge 1.5\,\mathrm{m}$ 需买全票；
  \item 水在标准大气压下沸腾温度 $t = 100\,^\circ\mathrm{C}$，结冰时 $t \le 0\,^\circ\mathrm{C}$；
  \item 身高：篮球运动员身高 $h > 1.9\,\mathrm{m}$；
  \item 预算：本月支出 $x \le 5000$ 元；
  \item 三角形两边之和大于第三边：$a+b>c$。
\end{itemize}
\small（答案不唯一，凡涉及大小、高低、长短、范围的实际关系均可）
\end{tcolorbox}
\end{frame}

% ========== 第2题 ==========
\begin{frame}[t]{第2题：投资方案}
\textbf{2.} 某市环保局为增加城市的绿地面积，提出两个投资方案：方案 A 为一次性投资 $500$ 万元；方案 B 为第一年投资 $100$ 万元，以后每年投资 $10$ 万元。列出不等式表示“经过 $n$ 年之后，方案 B 的投入不少于方案 A 的投入”。
\end{frame}

\begin{frame}[t]{第2题 解答}
\begin{tcolorbox}[colback=orange!10,colframe=orange,title=\textbf{解}]
方案 A 经过 $n$ 年的累计投入：
\[
S_A = 500\ \text{万元}.
\]

方案 B 第 $1$ 年投 $100$ 万元，第 $2$ 年起每年投 $10$ 万元，经过 $n$ 年的累计投入：
\[
S_B = 100 + 10(n-1)\ \text{万元}.
\]

“方案 B 的投入不少于方案 A” 即 $S_B \ge S_A$：
\[
\alert{100 + 10(n-1) \ge 500,\quad n\in \mathbf{N}^*}.
\]
\small（化简得 $n\ge 41$，即第 $41$ 年起方案 B 累计投入不少于方案 A）
\end{tcolorbox}
\end{frame}

% ========== 第3题 ==========
\begin{frame}[t]{第3题：作差比较大小}
\textbf{3.} 比较下列各组中两个代数式的大小：

\vspace{0.5em}
\begin{enumerate}[(1)]
  \item $x^2+5x+6$ 与 $2x^2+5x+9$；
  \item $(x-3)^2$ 与 $(x-2)(x-4)$；
  \item 当 $x>1$ 时，$x^2$ 与 $x^2-x+1$；
  \item $x^2+y^2+1$ 与 $2(x+y-1)$。
\end{enumerate}
\end{frame}

\begin{frame}[t]{第3题 解答}
\begin{tcolorbox}[colback=orange!10,colframe=orange,title=\textbf{解}]
(1) 作差：$(x^2+5x+6)-(2x^2+5x+9)=-x^2-3=-(x^2+3)<0$，故
\[
\alert{x^2+5x+6 < 2x^2+5x+9}.
\]

(2) 作差：$(x-3)^2-(x-2)(x-4)=(x^2-6x+9)-(x^2-6x+8)=1>0$，故
\[
\alert{(x-3)^2 > (x-2)(x-4)}.
\]

(3) 作差：$x^2-(x^2-x+1)=x-1$。因为 $x>1$，所以 $x-1>0$，故
\[
\alert{x^2 > x^2-x+1}.
\]

(4) 作差：$x^2+y^2+1-2(x+y-1)=x^2-2x+1+y^2-2y+1+1=(x-1)^2+(y-1)^2+1>0$，故
\[
\alert{x^2+y^2+1 > 2(x+y-1)}.
\]
\end{tcolorbox}
\end{frame}

% ========== 第4题 ==========
\begin{frame}[t]{第4题：两位数问题}
\textbf{4.} 一个大于 $50$ 且小于 $60$ 的两位数，其个位数字比十位数字大 $2$。试用不等式表示上述关系，并求出这个两位数（用 $a$ 和 $b$ 分别表示这个两位数的十位数字和个位数字）。
\end{frame}

\begin{frame}[t]{第4题 解答}
\begin{tcolorbox}[colback=orange!10,colframe=orange,title=\textbf{解}]
十位数字为 $a$，个位数字为 $b$，则这个两位数为 $10a+b$。由题意：
\[
\begin{cases}
50 < 10a+b < 60,\\
b = a+2,\\
a\in\{1,2,\dots,9\},\ b\in\{0,1,\dots,9\}.
\end{cases}
\]

把 $b=a+2$ 代入：
\[
50 < 10a+(a+2) < 60 \;\Longrightarrow\; 48 < 11a < 58.
\]
所以 $\dfrac{48}{11}<a<\dfrac{58}{11}$，即 $4.36<a<5.27$。又 $a\in\mathbf{N}^*$，故 $a=5$，$b=7$。

\vspace{0.5em}
所求两位数为 $\alert{57}$。
\end{tcolorbox}
\end{frame}

% ========== 第5题 ==========
\begin{frame}[t]{第5题：求代数式的取值范围}
\textbf{5.} 已知 $2<a<3$，$-2<b<-1$，求 $2a+b$ 的取值范围。
\end{frame}

\begin{frame}[t]{第5题 解答}
\begin{tcolorbox}[colback=orange!10,colframe=orange,title=\textbf{解}]
由 $2<a<3$，两边同乘 $2$（正数，不等号方向不变）：
\[
4 < 2a < 6.
\]
又已知
\[
-2 < b < -1.
\]
两个同向不等式相加（性质5）：
\[
4+(-2) < 2a+b < 6+(-1),
\]
即
\[
\alert{2 < 2a+b < 5}.
\]
\end{tcolorbox}
\end{frame}

% ========== 第6题 ==========
\begin{frame}[t]{第6题：证明传递性}
\textbf{6.} 证明：$c<b$，$b<a \Rightarrow c<a$。
\end{frame}

\begin{frame}[t]{第6题 解答}
\begin{tcolorbox}[colback=orange!10,colframe=orange,title=\textbf{证明}]
由 $c<b$，根据两实数大小关系的基本事实：
\[
c-b<0.
\]
由 $b<a$：
\[
b-a<0.
\]
作差比较 $c$ 与 $a$：
\[
c-a = (c-b)+(b-a).
\]
两个负数之和仍为负数，故 $c-a<0$，即 $c<a$。

\vspace{0.5em}
因此 $c<b$，$b<a \Rightarrow \alert{c<a}$。
\end{tcolorbox}
\end{frame}

% ========== 第7题 ==========
\begin{frame}[t]{第7题：综合证明}
\textbf{7.} 已知 $a>b>0$，$c<d<0$，$e<0$，求证：
\[
\frac{e}{a-c} > \frac{e}{b-d}.
\]
\end{frame}

\begin{frame}[t]{第7题 解答}
\begin{tcolorbox}[colback=orange!10,colframe=orange,title=\textbf{证明}]
\textbf{第一步：} 由 $c<d<0$，两边同乘 $-1$，不等号反向：
\[
-c > -d > 0.
\]
又 $a>b>0$，两式同向相加：
\[
a+(-c) > b+(-d) > 0 \;\Longrightarrow\; a-c > b-d > 0.
\]

\textbf{第二步：} 由 $a-c>b-d>0$，两边取倒数（正数，不等号反向）：
\[
\frac{1}{a-c} < \frac{1}{b-d}.
\]

\textbf{第三步：} 两边同乘负数 $e<0$，不等号反向：
\[
\alert{\frac{e}{a-c} > \frac{e}{b-d}}.
\]
\end{tcolorbox}
\end{frame}

% ========== 第8题 ==========
\begin{frame}[t]{第8题：选择题}
\textbf{8.} 下列命题为真命题的是（\quad）。

\vspace{0.5em}
(A) 若 $a>b>0$，则 $ac^2>bc^2$；\hfill
(B) 若 $a>b>0$，则 $a^2>b^2$；

\vspace{0.3em}
(C) 若 $a<b<0$，则 $a^2<ab<b^2$；\hfill
(D) 若 $a<b<0$，则 $\dfrac{1}{a}<\dfrac{1}{b}$。
\end{frame}

\begin{frame}[t]{第8题 解答}
\begin{tcolorbox}[colback=orange!10,colframe=orange,title=\textbf{解}]
\textbf{(A)} 当 $c=0$ 时，$ac^2=bc^2=0$，不成立。假命题。

\textbf{(B)} 由 $a>b>0$，据性质7（乘方）：$a^2>b^2$。\alert{真命题}。

\textbf{(C)} 由 $a<b<0$，两边乘 $a<0$（反向）：$a^2>ab$；两边乘 $b<0$（反向）：$ab>b^2$。故 $a^2>ab>b^2$。假命题。

\textbf{(D)} 由 $a<b<0$，取倒数（负数，反向）：$\dfrac{1}{a}>\dfrac{1}{b}$（如 $-2<-1$，而 $-\tfrac12>-1$）。假命题。

\vspace{0.5em}
综上，选 \alert{(B)}。
\end{tcolorbox}
\end{frame}

% ========== 第9题 ==========
\begin{frame}[t]{第9题：圆与正方形的面积}
\textbf{9.} 证明：圆的面积大于与它具有相同周长的正方形的面积。并据此说明，人们通常把自来水管的横截面制成圆形，而不是正方形的原因。
\end{frame}

\begin{frame}[t]{第9题 解答}
\begin{tcolorbox}[colback=orange!10,colframe=orange,title=\textbf{证明}]
设周长均为 $C$。

\textbf{圆：} 半径 $r=\dfrac{C}{2\pi}$，面积
\[
S_{\text{圆}}=\pi r^2=\pi\cdot\frac{C^2}{4\pi^2}=\frac{C^2}{4\pi}.
\]

\textbf{正方形：} 边长 $a=\dfrac{C}{4}$，面积
\[
S_{\text{正}}=a^2=\frac{C^2}{16}.
\]

因为 $4\pi\approx 12.57<16$，所以 $\dfrac{1}{4\pi}>\dfrac{1}{16}$，故
\[
\alert{S_{\text{圆}}=\frac{C^2}{4\pi} > \frac{C^2}{16}=S_{\text{正}}}.
\]

\vspace{0.3em}
\small\textbf{实际意义：} 周长（材料用量）相同时，圆形截面面积最大，水的流量最大；或相同流量时圆形最省材料。故自来水管横截面制成圆形。
\end{tcolorbox}
\end{frame}

% ========== 第10题 ==========
\begin{frame}[t]{第10题：糖水不等式}
\textbf{10.} 已知 $b$ 克糖水中含有 $a$ 克糖（$b>a>0$），再添加 $m$ 克糖（$m>0$）（假设全部溶解），糖水变甜了。请将这一事实表示为一个不等式，并证明这个不等式成立。
\end{frame}

\begin{frame}[t]{第10题 解答}
\begin{tcolorbox}[colback=orange!10,colframe=orange,title=\textbf{解}]
原来糖水浓度为 $\dfrac{a}{b}$，加糖后浓度为 $\dfrac{a+m}{b+m}$。

“糖水变甜” 即浓度变大：
\[
\frac{a+m}{b+m} > \frac{a}{b}.
\]

\textbf{证明（作差）：}
\[
\frac{a+m}{b+m}-\frac{a}{b}
= \frac{b(a+m)-a(b+m)}{b(b+m)}
= \frac{m(b-a)}{b(b+m)}.
\]
因为 $b>a>0$，$m>0$，所以分子 $m(b-a)>0$，分母 $b(b+m)>0$，故差 $>0$。

\vspace{0.3em}
即 \alert{$\dfrac{a+m}{b+m}>\dfrac{a}{b}$}。\small（糖水不等式）
\end{tcolorbox}
\end{frame}

% ========== 第11题 ==========
\begin{frame}[t]{第11题：证明算术根不等式}
\textbf{11.} 已知 $a>b>0$，求证：$\sqrt{a}>\sqrt{b}$。
\end{frame}

\begin{frame}[t]{第11题 解答}
\begin{tcolorbox}[colback=orange!10,colframe=orange,title=\textbf{证明}]
\textbf{证法一（作差）：}
\[
\sqrt{a}-\sqrt{b}
= \frac{(\sqrt{a}-\sqrt{b})(\sqrt{a}+\sqrt{b})}{\sqrt{a}+\sqrt{b}}
= \frac{a-b}{\sqrt{a}+\sqrt{b}}.
\]
因为 $a>b>0$，所以 $a-b>0$，且 $\sqrt{a}+\sqrt{b}>0$，故 $\sqrt{a}-\sqrt{b}>0$，即 $\sqrt{a}>\sqrt{b}$。

\vspace{0.5em}
\textbf{证法二（反证/性质）：} 假设 $\sqrt{a}\le\sqrt{b}$。因两边非负，据性质7平方得 $a\le b$，与已知 $a>b$ 矛盾。故 $\sqrt{a}>\sqrt{b}$。
\end{tcolorbox}
\end{frame}

% ========== 第12题 ==========
\begin{frame}[t]{第12题：火车站运货方案}
\textbf{12.} 火车站有某公司待运的甲种货物 $1530\,\mathrm{t}$，乙种货物 $1150\,\mathrm{t}$。现计划用 A、B 两种型号的货厢共 $50$ 节运送这批货物。已知 $35\,\mathrm{t}$ 甲种货物和 $15\,\mathrm{t}$ 乙种货物可装满一节 A 型货厢，$25\,\mathrm{t}$ 甲种货物和 $35\,\mathrm{t}$ 乙种货物可装满一节 B 型货厢，据此安排 A、B 两种货厢的节数，共有几种方案？若每节 A 型货厢的运费是 $0.5$ 万元，每节 B 型货厢的运费是 $0.8$ 万元，哪种方案的运费较少？
\end{frame}

\begin{frame}[t]{第12题 解答}
\begin{tcolorbox}[colback=orange!10,colframe=orange,title=\textbf{解}]
设安排 A 型货厢 $x$ 节，则 B 型货厢 $(50-x)$ 节。由题意：
\[
\begin{cases}
35x+25(50-x) \ge 1530, & \text{（甲货物）}\\
15x+35(50-x) \ge 1150. & \text{（乙货物）}
\end{cases}
\]

化简：
\[
\begin{cases}
10x \ge 280 \;\Rightarrow\; x\ge 28,\\
-20x \ge -600 \;\Rightarrow\; x\le 30.
\end{cases}
\]
又 $x\in\mathbf{N}^*$，故 $x=28,29,30$，共 \alert{3 种方案}：
\begin{itemize}
  \item A $28$ 节、B $22$ 节；
  \item A $29$ 节、B $21$ 节；
  \item A $30$ 节、B $20$ 节。
\end{itemize}

总运费 $y=0.5x+0.8(50-x)=40-0.3x$（万元）。$x$ 越大 $y$ 越小，故 $x=30$ 时运费最少：
\[
y_{\min}=40-0.3\times 30=\alert{31\ \text{万元}}.
\]
即安排 A 型 $30$ 节、B 型 $20$ 节时运费最少。
\end{tcolorbox}
\end{frame}

\begin{frame}[c]
  \begin{center}
    {\Huge\calligra Thanks!}
  \end{center}
\end{frame}

\end{document}
